\Needspace{14\baselineskip}

\Needspace{27\baselineskip}

\CNsection{6.1}{E-UTRAN Architecture}

\CNchained\CNsubsection{}{Why No RNC? The Flat Architecture Decision}

In UMTS (3G), the Radio Network Controller (RNC) sat between the NodeBs and the Core Network, centralizing radio resource management, handover decisions, and protocol termination. LTE \textbf{eliminated the RNC entirely}, creating a flat, single-node radio access network.

\textbf{Reasons for removing the RNC:}

{\def\LTcaptype{none} % do not increment counter
\Needspace{14\baselineskip}
\begin{longtable}[c]{@{}
  >{\raggedright\arraybackslash\hspace{0pt}}m{(\linewidth - 4\tabcolsep) * \real{0.2000}}
  >{\raggedright\arraybackslash\hspace{0pt}}m{(\linewidth - 4\tabcolsep) * \real{0.3333}}
  >{\raggedright\arraybackslash\hspace{0pt}}m{(\linewidth - 4\tabcolsep) * \real{0.4667}}@{}}
\toprule\noalign{}
\rowcolor{CNink}\CNth{Concern} & \CNth{3G (with RNC)} & \CNth{LTE (flat, no RNC)} \\
\CNheadrulesep
\endhead
\bottomrule\noalign{}
\endlastfoot
Latency & Extra hop through RNC adds \textasciitilde{}5–10~ms & Direct eNB\ensuremath{\leftrightarrow}EPC path, sub-5~ms user-plane latency \\
Single point of failure & RNC failure affects all connected NodeBs & Each eNB is independent \\
Scalability & RNC capacity limits number of NodeBs & No bottleneck; eNBs scale independently \\
Complexity & Three-node RAN (NodeB\ensuremath{\to}RNC\ensuremath{\to}CN) & Two-node RAN (eNB\ensuremath{\to}EPC) \\
Cost & Expensive, proprietary RNC hardware & Distributed intelligence reduces CAPEX \\
Handover speed & RNC-mediated = slower & Direct eNB-to-eNB via X2 = faster \\
\end{longtable}
}

\CNkeepfig{m06-enb-epc}{2}

\textbf{Key architectural principle:} All radio-related intelligence is pushed into the \textbf{eNodeB (eNB)}, which communicates directly with the Evolved Packet Core (EPC) via the S1 interface and with neighboring eNBs via the X2 interface.

\CNfigure{m06-enb-epc}{Each eNB connects to the MME over S1-MME and to the S-GW over S1-U}

\CNsection{6.2}{eNodeB Functions}

The eNodeB is the \textbf{sole node} in E-UTRAN. It performs all functions previously split between NodeB and RNC:

\CNsubsection{}{Radio Resource Management (RRM)}

\begin{itemize}
\tightlist
\item
  \textbf{Radio Bearer Control:} Establishment, maintenance, and release of radio bearers
\item
  \textbf{Radio Admission Control (RAC):} Accept/\allowbreak{}reject new radio bearer requests based on interference and load
\item
  \textbf{Connection Mobility Control:} Handover decisions based on UE measurements
\item
  \textbf{Dynamic Resource Allocation (Scheduling):} Assigns resource blocks to UEs every TTI (1~ms)
\end{itemize}

\CNsubsection{}{Scheduling}

\begin{itemize}
\tightlist
\item
  Operates every \textbf{1~ms TTI} (Transmission Time Interval)
\item
  Frequency-domain and time-domain scheduling
\item
  Algorithms: Round Robin, Proportional Fair, Max Throughput
\item
  Uses \textbf{CQI} (Channel Quality Indicator) reports from UEs
\item
  Controls both DL and UL resource allocation
\end{itemize}

\CNsubsection{}{HARQ (Hybrid Automatic Repeat Request)}

\begin{itemize}
\tightlist
\item
  Stop-and-wait HARQ with \textbf{8 parallel processes}
\item
  Synchronous in UL, asynchronous in DL
\item
  Soft combining: Chase Combining or Incremental Redundancy
\item
  Round-trip time: 8~ms (4~ms each way)
\end{itemize}

\CNsubsection{}{Channel Encoding}

\begin{itemize}
\tightlist
\item
  \textbf{Turbo coding} (rate 1/\allowbreak{}3) for data channels
\item
  \textbf{Convolutional coding} for control channels
\item
  Rate matching and HARQ buffer management
\item
  CRC attachment for error detection
\end{itemize}

\CNsubsection{}{Other Functions}

\begin{itemize}
\tightlist
\item
  IP header compression (ROHC)
\item
  Ciphering and integrity protection
\item
  Paging message distribution
\item
  Broadcast information (SIBs, MIBs)
\item
  Measurement configuration and reporting
\end{itemize}

\Needspace{14\baselineskip}

\CNkeepfig{m06-protocol-stack}{8}

\CNsection{6.3}{LTE Protocol Stack}

\CNchained\CNsubsection{}{Protocol Stack Diagram}

\CNfigure{m06-protocol-stack}{LTE protocol stacks of the UE, the eNodeB and the EPC}

\CNkeepfig{m06-cp-stack}{3}

\CNsubsection{}{Control Plane Stack}

\CNfigure{m06-cp-stack}{LTE control-plane protocol stack}

\CNkeepfig{m06-up-stack}{3}

\CNsubsection{}{User Plane Stack}

\CNfigure{m06-up-stack}{LTE user-plane protocol stack}

\Needspace{16\baselineskip}

\CNsection{6.4}{Protocol Layers in Detail}

\CNchained\CNsubsection{6.4.1}{Physical Layer (PHY)}

\CNchained\CNsubsubsection{Downlink: OFDMA (Orthogonal Frequency Division Multiple Access)}

\begin{itemize}
\tightlist
\item
  Divides bandwidth into many narrow \textbf{subcarriers} (15~kHz spacing)
\item
  Each subcarrier is orthogonal \ensuremath{\to} no inter-carrier interference
\item
  Robust against multipath fading (cyclic prefix absorbs delay spread)
\item
  Enables flexible frequency-domain scheduling
\end{itemize}

\textbf{Resource Grid:}

\begin{itemize}
\tightlist
\item
  \textbf{Resource Element (RE):} 1 subcarrier \ensuremath{\times} 1 OFDM symbol = smallest allocatable unit
\item
  \textbf{Resource Block (RB):} 12 subcarriers \ensuremath{\times} 7 OFDM symbols (1 slot = 0.5~ms) = \textbf{84~REs}
\item
  \textbf{Resource Block Pair:} 2~RBs = 1 subframe (1~ms TTI)
\end{itemize}

{\def\LTcaptype{none} % do not increment counter
\Needspace{14\baselineskip}
\begin{longtable}[c]{@{}ccc@{}}
\toprule\noalign{}
\rowcolor{CNink}\CNth{Bandwidth} & \CNth{Number of RBs} & \CNth{Subcarriers} \\
\CNheadrulesep
\endhead
\bottomrule\noalign{}
\endlastfoot
1.4~MHz & 6 & 72 \\
3~MHz & 15 & 180 \\
5~MHz & 25 & 300 \\
10~MHz & 50 & 600 \\
15~MHz & 75 & 900 \\
20~MHz & 100 & 1200 \\
\end{longtable}
}

\textbf{Frame Structure (FDD):}

\begin{itemize}
\tightlist
\item
  1 radio frame = 10~ms = 10 subframes
\item
  1 subframe = 1~ms = 2 slots
\item
  1 slot = 0.5~ms = 7 OFDM symbols (normal CP)
\end{itemize}

\CNsubsubsection{Uplink: SC-FDMA (Single Carrier FDMA)}

\begin{itemize}
\tightlist
\item
  Also called DFT-spread OFDM (DFT-s-OFDM)
\item
  Lower \textbf{PAPR} (Peak-to-Average Power Ratio) than OFDMA \ensuremath{\to} better for battery-powered UEs
\item
  Single-carrier property achieved by DFT precoding before IFFT
\item
  Same resource grid structure as DL (12 subcarriers per RB)
\item
  \textbf{Constraint:} UL allocation must be contiguous in frequency
\end{itemize}

\CNsubsubsection{Key PHY Procedures}

\begin{itemize}
\tightlist
\item
  Cell search (PSS/\allowbreak{}SSS synchronization)
\item
  Channel estimation (reference signals: CRS, DM-RS, SRS)
\item
  MIMO: up to 4\ensuremath{\times}4 in DL, 1\ensuremath{\times}2 or 2\ensuremath{\times}2 in UL (Rel-8/\allowbreak{}10)
\item
  Link adaptation via CQI/\allowbreak{}PMI/\allowbreak{}RI feedback
\item
  Power control (open-loop and closed-loop)
\end{itemize}

\CNsubsection{6.4.2}{MAC Layer (Medium Access Control)}

The MAC layer sits between RLC and PHY, responsible for \textbf{real-time scheduling} and data multiplexing.

\CNsubsubsection{Key Functions}

\begin{enumerate}
\def\labelenumi{\arabic{enumi}.}
\tightlist
\item
  \textbf{Scheduling (DL and UL)}

  \begin{itemize}
  \tightlist
  \item
    Operates every 1~ms TTI
  \item
    Assigns RBs based on CQI, buffer status, QoS priority, fairness
  \item
    DL: eNB decides; UL: eNB grants resources to UEs
  \end{itemize}
\item
  \textbf{Multiplexing/\allowbreak{}Demultiplexing}

  \begin{itemize}
  \tightlist
  \item
    Multiplexes data from multiple logical channels into one transport block
  \item
    Uses Logical Channel Prioritization (LCP) in UL
  \end{itemize}
\item
  \textbf{HARQ Entity}

  \begin{itemize}
  \tightlist
  \item
    8 parallel stop-and-wait processes
  \item
    DL: Asynchronous adaptive HARQ
  \item
    UL: Synchronous HARQ (fixed timing)
  \item
    Generates ACK/\allowbreak{}NACK; soft combining at receiver
  \end{itemize}
\item
  \textbf{Random Access Procedure}

  \begin{itemize}
  \tightlist
  \item
    Contention-based (4-step) or contention-free (2-step)
  \item
    Used for initial access, handover, UL synchronization recovery
  \end{itemize}
\item
  \textbf{BSR (Buffer Status Report)}

  \begin{itemize}
  \tightlist
  \item
    UE informs eNB about pending UL data
  \item
    Triggers UL scheduling grants
  \end{itemize}
\item
  \textbf{PHR (Power Headroom Report)}

  \begin{itemize}
  \tightlist
  \item
    UE reports remaining power margin
  \item
    Helps eNB with UL scheduling decisions
  \end{itemize}
\end{enumerate}

\CNkeepfig{m06-channel-levels}{2}

\CNsubsubsection{Channel Mapping at MAC}

\CNfigure{m06-channel-levels}{Three levels of LTE channels}

\Needspace{17\baselineskip}

\CNsubsection{6.4.3}{RLC Layer (Radio Link Control)}

RLC provides segmentation, reassembly, and (optionally) ARQ.

\CNsubsubsection{Three Modes}

{\def\LTcaptype{none} % do not increment counter
\Needspace{9\baselineskip}
\begin{longtable}[c]{@{}
  >{\raggedright\arraybackslash\hspace{0pt}}m{(\linewidth - 8\tabcolsep) * \real{0.1304}}
  >{\raggedright\arraybackslash\hspace{0pt}}m{(\linewidth - 8\tabcolsep) * \real{0.2391}}
  >{\raggedright\arraybackslash\hspace{0pt}}m{(\linewidth - 8\tabcolsep) * \real{0.1087}}
  >{\raggedright\arraybackslash\hspace{0pt}}m{(\linewidth - 8\tabcolsep) * \real{0.3043}}
  >{\raggedright\arraybackslash\hspace{0pt}}m{(\linewidth - 8\tabcolsep) * \real{0.2174}}@{}}
\toprule\noalign{}
\rowcolor{CNink}\CNth{Mode} & \CNth{Full Name} & \CNth{ARQ} & \CNth{Segmentation} & \CNth{Use Case} \\
\CNheadrulesep
\endhead
\bottomrule\noalign{}
\endlastfoot
\textbf{TM} & Transparent Mode & No & No & BCCH, PCCH, CCCH (small, fixed-size messages) \\
\textbf{UM} & Unacknowledged Mode & No & Yes & VoIP, real-time video (delay-sensitive) \\
\textbf{AM} & Acknowledged Mode & Yes & Yes & TCP traffic, signaling (reliability needed) \\
\end{longtable}
}

\CNsubsubsection{Key Functions}

\begin{itemize}
\tightlist
\item
  \textbf{Segmentation and Reassembly:} Splits PDCP PDUs into RLC-sized segments fitting transport block size
\item
  \textbf{Concatenation:} Packs multiple small SDUs into one PDU
\item
  \textbf{ARQ (AM only):} Retransmits lost RLC PDUs based on STATUS reports

  \begin{itemize}
  \tightlist
  \item
    Polling mechanism: RLC sender requests STATUS from receiver
  \item
    Operates independently from HARQ (RLC ARQ catches what HARQ misses)
  \end{itemize}
\item
  \textbf{Reordering (UM/\allowbreak{}AM):} Delivers SDUs in order to upper layers
\item
  \textbf{Duplicate Detection (AM):} Discards duplicate PDUs
\end{itemize}

\CNsubsubsection{RLC vs. HARQ}

\begin{itemize}
\tightlist
\item
  HARQ: Fast (8~ms RTT), limited retransmissions (typically 4)
\item
  RLC ARQ: Slower, but catches residual errors after HARQ exhaustion
\item
  Combined: achieves target BLER of 10\textsuperscript{-6} for AM bearers
\end{itemize}

\CNsubsection{6.4.4}{PDCP Layer (Packet Data Convergence Protocol)}

PDCP sits above RLC, providing security and efficiency functions.

\CNsubsubsection{Key Functions}

\begin{enumerate}
\def\labelenumi{\arabic{enumi}.}
\tightlist
\item
  \textbf{Header Compression (User Plane only)}

  \begin{itemize}
  \tightlist
  \item
    Uses \textbf{ROHC} (Robust Header Compression, RFC 3095/\allowbreak{}4995)
  \item
    Compresses IP/\allowbreak{}UDP/\allowbreak{}RTP headers from \textasciitilde{}40–60~bytes \ensuremath{\to} 1–4~bytes
  \item
    Critical for VoIP: 40-byte header on 32-byte payload \ensuremath{\to} \textasciitilde{}56\% overhead without ROHC
  \item
    Multiple ROHC profiles for different protocol stacks
  \end{itemize}
\item
  \textbf{Ciphering (User Plane + Control Plane)}

  \begin{itemize}
  \tightlist
  \item
    Encrypts all user data and signaling
  \item
    Algorithms: EEA0 (null), 128-EEA1 (SNOW 3G), 128-EEA2 (AES-CTR), 128-EEA3 (ZUC)
  \item
    Uses COUNT (sequence number + HFN) + BEARER + DIRECTION as input
  \end{itemize}
\item
  \textbf{Integrity Protection (Control Plane only)}

  \begin{itemize}
  \tightlist
  \item
    Protects RRC signaling from tampering
  \item
    Algorithms: EIA1 (SNOW 3G), EIA2 (AES-CMAC), EIA3 (ZUC)
  \item
    Generates 32-bit MAC-I appended to each RRC message
  \item
    \textbf{Not applied to user plane in LTE} (added in 5G NR for user plane)
  \end{itemize}
\item
  \textbf{Sequence Numbering and Reordering}

  \begin{itemize}
  \tightlist
  \item
    Assigns PDCP SN to each SDU
  \item
    Reorders out-of-sequence PDUs (important during handover)
  \item
    Detects duplicates
  \end{itemize}
\item
  \textbf{Handover Support}

  \begin{itemize}
  \tightlist
  \item
    In-sequence delivery and duplicate elimination during handover
  \item
    Status reporting: tells source eNB which PDUs were successfully delivered
  \item
    Enables lossless handover for AM bearers
  \end{itemize}
\end{enumerate}

\CNkeepfig{m06-pdcp-pdu}{2}

\CNsubsubsection{PDCP PDU Format}

\CNfigure{m06-pdcp-pdu}{PDCP PDU format}

\CNsubsection{6.4.5}{RRC Layer (Radio Resource Control)}

RRC is the \textbf{control plane protocol} between UE and eNodeB, managing the radio connection.

\CNsubsubsection{Key Functions}

\begin{enumerate}
\def\labelenumi{\arabic{enumi}.}
\tightlist
\item
  \textbf{Connection Management}

  \begin{itemize}
  \tightlist
  \item
    RRC Connection Setup /\allowbreak{} Reconfiguration /\allowbreak{} Release
  \item
    Security activation (ciphering + integrity)
  \item
    SRB (Signaling Radio Bearer) establishment
  \end{itemize}
\item
  \textbf{Measurement Configuration and Reporting}

  \begin{itemize}
  \tightlist
  \item
    Configures UE to measure serving and neighbor cells
  \item
    Events: A1–A6 (intra/\allowbreak{}inter-frequency), B1–B2 (inter-RAT)
  \item
    Reporting: periodic or event-triggered
  \item
    Quantities: RSRP, RSRQ, SINR
  \end{itemize}
\item
  \textbf{Handover Control}

  \begin{itemize}
  \tightlist
  \item
    eNB makes handover decision based on measurement reports
  \item
    Sends \CNcode{RRCConnectionReconfiguration} with mobility control info
  \item
    UE performs handover, sends \CNcode{RRCConnectionReconfigurationComplete}
  \end{itemize}
\item
  \textbf{System Information Broadcasting}

  \begin{itemize}
  \tightlist
  \item
    \textbf{MIB:} Transmitted on PBCH every 40~ms (bandwidth, SFN, PHICH config)
  \item
    \textbf{SIB1:} Cell access info, scheduling of other SIBs (80~ms period)
  \item
    \textbf{SIB2:} Common channel config (RACH, paging, UL power control)
  \item
    \textbf{SIB3–SIB8:} Cell reselection, inter-frequency/\allowbreak{}inter-RAT info
  \item
    \textbf{SIB9–SIB19:} Additional info (home eNB, MBMS, etc.)
  \end{itemize}
\item
  \textbf{Paging}

  \begin{itemize}
  \tightlist
  \item
    Notifies UEs in IDLE of incoming calls/\allowbreak{}data
  \item
    DRX-based: UE wakes at specific paging occasions
  \end{itemize}
\item
  \textbf{NAS Message Transfer}

  \begin{itemize}
  \tightlist
  \item
    Transparently carries NAS messages between UE and MME
  \end{itemize}
\end{enumerate}

\Needspace{12\baselineskip}

\CNsubsubsection{Signaling Radio Bearers (SRBs)}

{\def\LTcaptype{none} % do not increment counter
\Needspace{9\baselineskip}
\begin{longtable}[c]{@{}ccc@{}}
\toprule\noalign{}
\rowcolor{CNink}\CNth{SRB} & \CNth{Purpose} & \CNth{RLC Mode} \\
\CNheadrulesep
\endhead
\bottomrule\noalign{}
\endlastfoot
SRB0 & RRC messages on CCCH (before security) & TM \\
SRB1 & RRC messages + piggybacked NAS (after setup) & AM \\
SRB2 & NAS messages (lower priority than SRB1) & AM \\
\end{longtable}
}

\Needspace{14\baselineskip}

\Needspace{20\baselineskip}

\CNsection{6.5}{LTE Channel Structure}

\CNchained\CNsubsection{}{Channel Types Overview}

LTE defines three categories of channels, each at a different layer of abstraction:

{\def\LTcaptype{none} % do not increment counter
\Needspace{9\baselineskip}
\begin{longtable}[c]{@{}
  >{\raggedright\arraybackslash\hspace{0pt}}m{(\linewidth - 6\tabcolsep) * \real{0.2128}}
  >{\raggedright\arraybackslash\hspace{0pt}}m{(\linewidth - 6\tabcolsep) * \real{0.1489}}
  >{\raggedright\arraybackslash\hspace{0pt}}m{(\linewidth - 6\tabcolsep) * \real{0.2340}}
  >{\raggedright\arraybackslash\hspace{0pt}}m{(\linewidth - 6\tabcolsep) * \real{0.4043}}@{}}
\toprule\noalign{}
\rowcolor{CNink}\CNth{Category} & \CNth{Layer} & \CNth{Defined By} & \CNth{Question Answered} \\
\CNheadrulesep
\endhead
\bottomrule\noalign{}
\endlastfoot
Logical Channels & MAC (upper) & \textbf{What} type of data & What is being transferred? \\
Transport Channels & MAC (lower) & \textbf{How} data is transferred & How is it characterized? \\
Physical Channels & PHY & \textbf{Where} in the resource grid & Where are the bits placed? \\
\end{longtable}
}

\Needspace{11\baselineskip}

\Needspace{19\baselineskip}

\CNsubsection{}{Logical Channels}

\CNchained\CNsubsubsection{Control Channels}

{\def\LTcaptype{none} % do not increment counter
\Needspace{13\baselineskip}
\begin{longtable}[c]{@{}
  >{\raggedright\arraybackslash\hspace{0pt}}m{(\linewidth - 6\tabcolsep) * \real{0.2571}}
  >{\raggedright\arraybackslash\hspace{0pt}}m{(\linewidth - 6\tabcolsep) * \real{0.1714}}
  >{\raggedright\arraybackslash\hspace{0pt}}m{(\linewidth - 6\tabcolsep) * \real{0.3143}}
  >{\raggedright\arraybackslash\hspace{0pt}}m{(\linewidth - 6\tabcolsep) * \real{0.2571}}@{}}
\toprule\noalign{}
\rowcolor{CNink}\CNth{Channel} & \CNth{Name} & \CNth{Direction} & \CNth{Purpose} \\
\CNheadrulesep
\endhead
\bottomrule\noalign{}
\endlastfoot
\textbf{BCCH} & Broadcast Control Channel & DL & System information (MIB, SIBs) \\
\textbf{PCCH} & Paging Control Channel & DL & Paging messages for IDLE UEs \\
\textbf{CCCH} & Common Control Channel & DL/\allowbreak{}UL & Messages before RRC connection (Msg3/\allowbreak{}Msg4) \\
\textbf{DCCH} & Dedicated Control Channel & DL/\allowbreak{}UL & Dedicated signaling (RRC after connection) \\
\textbf{MCCH} & Multicast Control Channel & DL & MBMS control information \\
\end{longtable}
}

\Needspace{11\baselineskip}

\CNsubsubsection{Traffic Channels}

{\def\LTcaptype{none} % do not increment counter
\Needspace{8\baselineskip}
\begin{longtable}[c]{@{}
  >{\raggedright\arraybackslash\hspace{0pt}}m{(\linewidth - 6\tabcolsep) * \real{0.2571}}
  >{\raggedright\arraybackslash\hspace{0pt}}m{(\linewidth - 6\tabcolsep) * \real{0.1714}}
  >{\raggedright\arraybackslash\hspace{0pt}}m{(\linewidth - 6\tabcolsep) * \real{0.3143}}
  >{\raggedright\arraybackslash\hspace{0pt}}m{(\linewidth - 6\tabcolsep) * \real{0.2571}}@{}}
\toprule\noalign{}
\rowcolor{CNink}\CNth{Channel} & \CNth{Name} & \CNth{Direction} & \CNth{Purpose} \\
\CNheadrulesep
\endhead
\bottomrule\noalign{}
\endlastfoot
\textbf{DTCH} & Dedicated Traffic Channel & DL/\allowbreak{}UL & User data (IP packets) \\
\textbf{MTCH} & Multicast Traffic Channel & DL & MBMS user data \\
\end{longtable}
}

\Needspace{18\baselineskip}

\CNsubsection{}{Transport Channels}

{\def\LTcaptype{none} % do not increment counter
\Needspace{14\baselineskip}
\begin{longtable}[c]{@{}
  >{\raggedright\arraybackslash\hspace{0pt}}m{(\linewidth - 6\tabcolsep) * \real{0.2571}}
  >{\raggedright\arraybackslash\hspace{0pt}}m{(\linewidth - 6\tabcolsep) * \real{0.1714}}
  >{\raggedright\arraybackslash\hspace{0pt}}m{(\linewidth - 6\tabcolsep) * \real{0.3143}}
  >{\raggedright\arraybackslash\hspace{0pt}}m{(\linewidth - 6\tabcolsep) * \real{0.2571}}@{}}
\toprule\noalign{}
\rowcolor{CNink}\CNth{Channel} & \CNth{Name} & \CNth{Direction} & \CNth{Purpose} \\
\CNheadrulesep
\endhead
\bottomrule\noalign{}
\endlastfoot
\textbf{BCH} & Broadcast Channel & DL & Carries MIB; fixed format, 40~ms period \\
\textbf{PCH} & Paging Channel & DL & Carries paging; supports DRX \\
\textbf{DL-SCH} & Downlink Shared Channel & DL & Main DL data; supports HARQ, dynamic scheduling \\
\textbf{UL-SCH} & Uplink Shared Channel & UL & Main UL data; supports HARQ, dynamic scheduling \\
\textbf{RACH} & Random Access Channel & UL & Random access preambles \\
\textbf{MCH} & Multicast Channel & DL & MBMS data; fixed scheduling \\
\end{longtable}
}

\Needspace{11\baselineskip}

\Needspace{20\baselineskip}

\CNsubsection{}{Physical Channels}

\CNchained\CNsubsubsection{Downlink Physical Channels}

{\def\LTcaptype{none} % do not increment counter
\Needspace{14\baselineskip}
\begin{longtable}[c]{@{}
  >{\raggedright\arraybackslash\hspace{0pt}}m{(\linewidth - 4\tabcolsep) * \real{0.3750}}
  >{\raggedright\arraybackslash\hspace{0pt}}m{(\linewidth - 4\tabcolsep) * \real{0.2500}}
  >{\raggedright\arraybackslash\hspace{0pt}}m{(\linewidth - 4\tabcolsep) * \real{0.3750}}@{}}
\toprule\noalign{}
\rowcolor{CNink}\CNth{Channel} & \CNth{Name} & \CNth{Purpose} \\
\CNheadrulesep
\endhead
\bottomrule\noalign{}
\endlastfoot
\textbf{PDSCH} & Physical DL Shared Channel & User data + system info (SIBs, paging) \\
\textbf{PDCCH} & Physical DL Control Channel & DCI (scheduling grants, power control) \\
\textbf{PBCH} & Physical Broadcast Channel & MIB transmission \\
\textbf{PMCH} & Physical Multicast Channel & MBMS data \\
\textbf{PCFICH} & Physical CFI Channel & Indicates \# of OFDM symbols for PDCCH \\
\textbf{PHICH} & Physical HARQ Indicator Channel & UL HARQ ACK/\allowbreak{}NACK \\
\end{longtable}
}

\Needspace{12\baselineskip}

\CNsubsubsection{Uplink Physical Channels}

{\def\LTcaptype{none} % do not increment counter
\Needspace{9\baselineskip}
\begin{longtable}[c]{@{}ccc@{}}
\toprule\noalign{}
\rowcolor{CNink}\CNth{Channel} & \CNth{Name} & \CNth{Purpose} \\
\CNheadrulesep
\endhead
\bottomrule\noalign{}
\endlastfoot
\textbf{PUSCH} & Physical UL Shared Channel & UL user data + control info \\
\textbf{PUCCH} & Physical UL Control Channel & CQI, HARQ ACK/\allowbreak{}NACK, SR \\
\textbf{PRACH} & Physical Random Access Channel & Random access preambles \\
\end{longtable}
}

\CNkeepfig{m06-channel-mapping}{3}

\CNsubsection{}{Channel Mapping}

\CNfigure{m06-channel-mapping}{Mapping of logical, transport and physical channels}

\Needspace{26\baselineskip}

\CNsubsection{}{Complete Channel Mapping Table}

{\def\LTcaptype{none} % do not increment counter
\Needspace{22\baselineskip}
\begin{longtable}[c]{@{}cccc@{}}
\toprule\noalign{}
\rowcolor{CNink}\CNth{Logical Channel} & \CNth{Transport Channel} & \CNth{Physical Channel} & \CNth{Direction} \\
\CNheadrulesep
\endhead
\bottomrule\noalign{}
\endlastfoot
BCCH (MIB) & BCH & PBCH & DL \\
BCCH (SIBs) & DL-SCH & PDSCH & DL \\
PCCH & PCH & PDSCH & DL \\
CCCH & DL-SCH & PDSCH & DL \\
CCCH & UL-SCH & PUSCH & UL \\
DCCH & DL-SCH & PDSCH & DL \\
DCCH & UL-SCH & PUSCH & UL \\
DTCH & DL-SCH & PDSCH & DL \\
DTCH & UL-SCH & PUSCH & UL \\
— (HARQ feedback) & — & PHICH & DL \\
— (DCI/\allowbreak{}grants) & — & PDCCH & DL \\
— (CFI) & — & PCFICH & DL \\
— (CQI/\allowbreak{}SR/\allowbreak{}ACK) & — & PUCCH & UL \\
— (RA preamble) & RACH & PRACH & UL \\
\end{longtable}
}

\CNsection{6.6}{S1 Interface}

The S1 interface connects the eNodeB to the EPC. It is split into:

\begin{itemize}
\tightlist
\item
  \textbf{S1-MME (S1-C):} Control plane \ensuremath{\to} eNB \ensuremath{\leftrightarrow} MME
\item
  \textbf{S1-U:} User plane \ensuremath{\to} eNB \ensuremath{\leftrightarrow} S-GW
\end{itemize}

\CNkeepfig{m06-s1mme-stack}{3}

\CNsubsection{}{S1-MME Protocol Stack}

\CNfigure{m06-s1mme-stack}{S1-MME protocol stack}

\textbf{Why SCTP (not TCP)?}

\begin{itemize}
\tightlist
\item
  Multi-streaming: avoids head-of-line blocking
\item
  Multi-homing: failover between IP addresses
\item
  Message-oriented (preserves boundaries)
\end{itemize}

\Needspace{11\baselineskip}

\Needspace{26\baselineskip}

\CNsubsection{}{S1AP Procedures}

\CNchained\CNsubsubsection{UE-Associated Procedures (per-UE context)}

{\def\LTcaptype{none} % do not increment counter
\Needspace{20\baselineskip}
\begin{longtable}[c]{@{}
  >{\raggedright\arraybackslash\hspace{0pt}}m{(\linewidth - 4\tabcolsep) * \real{0.3548}}
  >{\raggedright\arraybackslash\hspace{0pt}}m{(\linewidth - 4\tabcolsep) * \real{0.3548}}
  >{\raggedright\arraybackslash\hspace{0pt}}m{(\linewidth - 4\tabcolsep) * \real{0.2903}}@{}}
\toprule\noalign{}
\rowcolor{CNink}\CNth{Procedure} & \CNth{Direction} & \CNth{Purpose} \\
\CNheadrulesep
\endhead
\bottomrule\noalign{}
\endlastfoot
Initial UE Message & eNB \ensuremath{\to} MME & Carries first NAS message (Attach, Service Request) \\
Downlink NAS Transport & MME \ensuremath{\to} eNB & Delivers NAS message to UE \\
Uplink NAS Transport & eNB \ensuremath{\to} MME & Forwards NAS message from UE \\
Initial Context Setup & MME \ensuremath{\to} eNB & Establishes security context + default bearer \\
UE Context Release & MME \ensuremath{\leftrightarrow} eNB & Releases UE context (Request/\allowbreak{}Command/\allowbreak{}Complete) \\
Handover Preparation & Source eNB \ensuremath{\to} MME & S1-based handover when X2 unavailable \\
Handover Resource Allocation & MME \ensuremath{\to} Target eNB & Allocates resources at target \\
Path Switch Request & Target eNB \ensuremath{\to} MME & Updates S-GW path after X2 handover \\
E-RAB Setup/\allowbreak{}Modify/\allowbreak{}Release & MME \ensuremath{\to} eNB & Bearer management \\
\end{longtable}
}

\Needspace{17\baselineskip}

\CNsubsubsection{Non-UE-Associated Procedures (global)}

{\def\LTcaptype{none} % do not increment counter
\Needspace{14\baselineskip}
\begin{longtable}[c]{@{}cc@{}}
\toprule\noalign{}
\rowcolor{CNink}\CNth{Procedure} & \CNth{Purpose} \\
\CNheadrulesep
\endhead
\bottomrule\noalign{}
\endlastfoot
S1 Setup & eNB registers with MME (TAC, PLMN, capacity) \\
Reset & Error recovery (partial or full) \\
Paging & MME sends paging to eNBs in tracking area \\
eNB/\allowbreak{}MME Configuration Update & Capacity/\allowbreak{}configuration changes \\
Overload Start/\allowbreak{}Stop & MME signals congestion \\
Error Indication & Report protocol errors \\
\end{longtable}
}

\CNkeepfig{m06-s1u-stack}{3}

\CNsubsection{}{S1-U Protocol Stack}

\CNfigure{m06-s1u-stack}{S1-U protocol stack}

\begin{itemize}
\tightlist
\item
  Each E-RAB has a \textbf{GTP tunnel} identified by TEID (Tunnel Endpoint ID)
\item
  One-to-one mapping: EPS bearer \ensuremath{\leftrightarrow} S1 GTP tunnel \ensuremath{\leftrightarrow} radio bearer
\end{itemize}

\CNsection{6.7}{X2 Interface}

The X2 interface connects \textbf{neighboring eNodeBs} directly for:

\begin{itemize}
\tightlist
\item
  Fast handover (no MME involvement for intra-LTE mobility)
\item
  Inter-cell load management
\item
  Inter-cell interference coordination (ICIC)
\end{itemize}

\CNkeepfig{m06-x2-stacks}{3}

\CNsubsection{}{X2 Protocol Stack}

\CNfigure{m06-x2-stacks}{X2 protocol stacks}

\Needspace{11\baselineskip}

\CNsubsection{}{X2AP Procedures}

\CNchained\CNsubsubsection{Mobility Management (Handover)}

\textbf{X2 Handover Sequence:}

\begin{enumerate}
\def\labelenumi{\arabic{enumi}.}
\tightlist
\item
  \textbf{Source eNB} sends \CNcode{Handover Request} \ensuremath{\to} Target eNB (includes UE context, bearer info)
\item
  \textbf{Target eNB} performs admission control \ensuremath{\to} sends \CNcode{Handover Request Acknowledge}
\item
  \textbf{Source eNB} sends \CNcode{RRCConnectionReconfiguration} to UE (with target cell info)
\item
  \textbf{UE} synchronizes to target cell, sends \CNcode{RRCConnectionReconfigurationComplete}
\item
  \textbf{Target eNB} sends \CNcode{Path Switch Request} to MME (updates core network path)
\item
  \textbf{Source eNB} receives \CNcode{UE Context Release} \ensuremath{\to} releases resources
\end{enumerate}

\textbf{During handover:} Data forwarded source\ensuremath{\to}target via X2-U GTP tunnel (avoids data loss).

\Needspace{12\baselineskip}

\CNsubsubsection{Load Management}

{\def\LTcaptype{none} % do not increment counter
\Needspace{9\baselineskip}
\begin{longtable}[c]{@{}
  >{\raggedright\arraybackslash\hspace{0pt}}m{(\linewidth - 2\tabcolsep) * \real{0.5500}}
  >{\raggedright\arraybackslash\hspace{0pt}}m{(\linewidth - 2\tabcolsep) * \real{0.4500}}@{}}
\toprule\noalign{}
\rowcolor{CNink}\CNth{Procedure} & \CNth{Purpose} \\
\CNheadrulesep
\endhead
\bottomrule\noalign{}
\endlastfoot
Resource Status Reporting & Request/\allowbreak{}report load information (PRB usage, composite load) \\
Load Indication & Inform neighbor about UL interference (HII, OI) \\
Mobility Change Request & Negotiate handover parameters (CIO adjustment) \\
\end{longtable}
}

\CNsubsubsection{Inter-Cell Interference Coordination (ICIC)}

\begin{itemize}
\tightlist
\item
  \textbf{Load Information message} carries:

  \begin{itemize}
  \tightlist
  \item
    \textbf{HII (High Interference Indication):} UL RBs where cell plans to schedule cell-edge UEs
  \item
    \textbf{OI (Overload Indicator):} UL RBs experiencing high interference
  \item
    \textbf{RNTP (Relative Narrowband TX Power):} DL RBs where cell will exceed power threshold
  \end{itemize}
\end{itemize}

\CNsubsection{}{X2 Setup and Configuration}

\begin{itemize}
\tightlist
\item
  X2 link established via \textbf{X2 Setup procedure} (or via MME using S1 ENB Configuration Transfer)
\item
  Automatic Neighbor Relation (ANR): eNB discovers neighbors and establishes X2
\item
  Multiple eNBs can connect forming a mesh topology
\end{itemize}

\Needspace{14\baselineskip}

\CNkeepfig{m06-rrc-states}{8}

\CNsection{6.8}{RRC States and Transitions}

\CNchained\CNsubsection{}{RRC State Machine}

\CNfigure{m06-rrc-states}{LTE RRC state machine}

\Needspace{20\baselineskip}

\CNsubsection{}{RRC\_\allowbreak{}IDLE State}

{\def\LTcaptype{none} % do not increment counter
\Needspace{16\baselineskip}
\begin{longtable}[c]{@{}cc@{}}
\toprule\noalign{}
\rowcolor{CNink}\CNth{Characteristic} & \CNth{Description} \\
\CNheadrulesep
\endhead
\bottomrule\noalign{}
\endlastfoot
Context & No eNB context; registered with MME only \\
Location & Known at \textbf{Tracking Area} level \\
Mobility & \textbf{Cell reselection} (UE-controlled, based on SIBs) \\
Data & No data transfer possible \\
Paging & UE monitors paging channel with DRX cycle \\
System Info & UE acquires MIB and relevant SIBs \\
Power & Low power consumption (long DRX sleep) \\
\end{longtable}
}

\Needspace{20\baselineskip}

\CNsubsection{}{RRC\_\allowbreak{}CONNECTED State}

{\def\LTcaptype{none} % do not increment counter
\Needspace{16\baselineskip}
\begin{longtable}[c]{@{}cc@{}}
\toprule\noalign{}
\rowcolor{CNink}\CNth{Characteristic} & \CNth{Description} \\
\CNheadrulesep
\endhead
\bottomrule\noalign{}
\endlastfoot
Context & Full UE context in serving eNB \\
Location & Known at \textbf{cell} level \\
Mobility & \textbf{Handover} (network-controlled, measurement-based) \\
Data & Active data transfer (DL + UL) \\
Measurements & Configured by eNB (events A1–A6, B1–B2) \\
DRX & Connected-mode DRX possible (shorter cycles) \\
Security & AS ciphering + integrity active \\
\end{longtable}
}

\Needspace{11\baselineskip}

\CNsubsection{}{State Transitions}

\CNchained\CNsubsubsection{IDLE \ensuremath{\to} CONNECTED (Connection Establishment)}

\begin{enumerate}
\def\labelenumi{\arabic{enumi}.}
\tightlist
\item
  UE sends \textbf{RRC Connection Request} (on CCCH via SRB0)
\item
  eNB responds with \textbf{RRC Connection Setup} (configures SRB1)
\item
  UE sends \textbf{RRC Connection Setup Complete} (carries initial NAS message)
\item
  eNB forwards NAS to MME via S1AP \CNcode{Initial UE Message}
\item
  MME initiates \textbf{Initial Context Setup} (security + bearer establishment)
\end{enumerate}

\CNsubsubsection{CONNECTED \ensuremath{\to} IDLE (Connection Release)}

\begin{itemize}
\tightlist
\item
  MME triggers via S1AP \CNcode{UE Context Release Command}
\item
  eNB sends \textbf{RRC Connection Release} to UE
\item
  UE deletes AS context, returns to idle
\item
  Triggers: inactivity timer, explicit NAS detach, radio link failure recovery failure
\end{itemize}

\CNkeepfig{m06-rrc-inactive}{5}

\CNsubsection{}{5G Extension: RRC\_\allowbreak{}INACTIVE State}

In 5G NR (and late LTE Rel-15+), a third state is introduced:

\CNfigure{m06-rrc-inactive}{5G NR RRC states with RRC\_\allowbreak{}INACTIVE}

{\def\LTcaptype{none} % do not increment counter
\Needspace{9\baselineskip}
\begin{longtable}[c]{@{}
  >{\raggedright\arraybackslash\hspace{0pt}}m{(\linewidth - 8\tabcolsep) * \real{0.1719}}
  >{\raggedright\arraybackslash\hspace{0pt}}m{(\linewidth - 8\tabcolsep) * \real{0.1719}}
  >{\raggedright\arraybackslash\hspace{0pt}}m{(\linewidth - 8\tabcolsep) * \real{0.2969}}
  >{\raggedright\arraybackslash\hspace{0pt}}m{(\linewidth - 8\tabcolsep) * \real{0.1562}}
  >{\raggedright\arraybackslash\hspace{0pt}}m{(\linewidth - 8\tabcolsep) * \real{0.2031}}@{}}
\toprule\noalign{}
\rowcolor{CNink}\CNth{State} & \CNth{UE Context} & \CNth{Location Granularity} & \CNth{Mobility} & \CNth{CN Connection} \\
\CNheadrulesep
\endhead
\bottomrule\noalign{}
\endlastfoot
IDLE & None in RAN & Tracking Area & Cell reselection & None \\
INACTIVE & Stored in gNB (suspended) & RAN Notification Area & Cell reselection + RAN paging & Maintained with CN \\
CONNECTED & Active in gNB & Cell level & Handover & Active \\
\end{longtable}
}

\textbf{Benefits of INACTIVE state:}

\begin{itemize}
\tightlist
\item
  Faster resumption than IDLE\ensuremath{\to}CONNECTED (skip NAS procedures)
\item
  Power saving (no connected-mode signaling)
\item
  Reduced signaling load on core network
\item
  AS context preserved \ensuremath{\to} faster resume with single RRC Resume message
\item
  Ideal for IoT devices with intermittent small data transmissions
\end{itemize}

\Needspace{28\baselineskip}

\CNsection{}{Summary}

{\def\LTcaptype{none} % do not increment counter
\Needspace{22\baselineskip}
\begin{longtable}[c]{@{}
  >{\raggedright\arraybackslash\hspace{0pt}}m{(\linewidth - 2\tabcolsep) * \real{0.4583}}
  >{\raggedright\arraybackslash\hspace{0pt}}m{(\linewidth - 2\tabcolsep) * \real{0.5417}}@{}}
\toprule\noalign{}
\rowcolor{CNink}\CNth{Component} & \CNth{Key Takeaway} \\
\CNheadrulesep
\endhead
\bottomrule\noalign{}
\endlastfoot
Architecture & Flat (no RNC); eNB handles all RAN functions \\
eNodeB & Scheduling, RRM, HARQ, handover decisions, encoding \\
PHY & OFDMA (DL) /\allowbreak{} SC-FDMA (UL); resource blocks as basic unit \\
MAC & 1~ms scheduling, HARQ, multiplexing, random access \\
RLC & TM/\allowbreak{}UM/\allowbreak{}AM; segmentation + ARQ for reliability \\
PDCP & ROHC compression, ciphering, integrity (CP), reordering \\
RRC & Connection/\allowbreak{}mobility management, measurement, system info \\
Channels & Logical\ensuremath{\to}Transport\ensuremath{\to}Physical mapping; shared channels dominate \\
S1 & eNB\ensuremath{\leftrightarrow}EPC; S1AP (SCTP) for control, GTP-U for user plane \\
X2 & eNB\ensuremath{\leftrightarrow}eNB; fast handover, load mgmt, interference coordination \\
RRC States & IDLE (low power) \ensuremath{\leftrightarrow} CONNECTED (active); INACTIVE in 5G \\
\end{longtable}
}

\emph{Module 6 — E-UTRAN (LTE Radio Access Network) — Advanced Communication Networks}
